\documentclass{report}
\usepackage{amsmath,amssymb}
\setlength{\parindent}{0mm}
\setlength{\parskip}{1em}
\begin{document}
\begin{center}
\rule{6in}{1pt} \
{\large Michal Kordy \\
{\bf Finite element solution of a Schelkunoff vector potential for frequency domain, EM field simulation}}

JWB 321 Department of Mathematics \\ University of Utah \\ 155 S 1400 E SALT LAKE CITY \\ UT 84112-0090
\\
{\tt mikordy@gmail.com}\\
Elena Cherkaev\\
Phil Wannamaker\end{center}

A novel method for the 3-D diffusive electromagnetic (EM) forward problem
is developed and tested. A Lorentz-gauge, Schelkunoff complex vector
potential is used to represent the EM field in the frequency domain and
the nodal finite element method is used for numerical simulation. The
potential allows for three degrees of freedom per node, instead of four
if Coulomb-gauge vector and scalar potentials are used. Unlike the
finite-difference method, which minimizes error at discrete points, the
finite element method minimizes error over the entire domain cell volumes
and may easily adapt to complex topography. Existence and uniqueness of
this continuous Schelkunoff potential is proven, boundary conditions are
found and a governing equation satisfied by the potential in weak form is
obtained. This approach for using a Schelkunoff potential in the finite
element method differs from others found in the literature. If the
standard weak form of the Helmholtz equation is used, the obtained
solution is continuous and has continuous normal derivative across
boundaries of regions with different physical properties; however,
continuous Schelkunoff potential components do not have continuous normal
derivative, divergence of the potential divided by (complex) conductivity
and magnetic permeability is continuous instead. Two weak forms of the
governing equation are tried. Both of them produce a system matrix that
is ill-conditioned and as a result iterative schemes do not converge. A
different idea is tried next, instead of representing electric field by a
Schelkunoff potential, magnetic field is represented in a similar way.
When it is assumed that magnetic permeability $\mu$ is constant, a
convenient weak form of the governing equation is obtained. This form is
tested numerically on a simple model of a conducting prism in a resistive
whole space and the code gives a similar results to an independent finite
difference code.


\end{document}
